% os-axioms.tex — Chapter on the Osterwalder–Schrader axioms for
% the continuum SU(2) Wilson measure mu_infinity.
%
% Linked code / Lean:
%   - lean/YMGap/Continuum/OSAxioms.lean (Lean type signature)
%   - \Cref{ch:os-axioms} of work_plan.md.

\chapter{Osterwalder--Schrader Axioms for the
  Continuum \texorpdfstring{$\SUtwo$}{SU(2)} Wilson Measure
  \texorpdfstring{$\muinf$}{mu\_infinity}}
\label{ch:os-axioms}

\noindent\emph{Verification of OS0--OS4 for $\muinf$, with $\sighinf$
delivering the explicit OS4 cluster rate.}

We verify the five Osterwalder--Schrader axioms (OS0--OS4) for the
continuum $\SUtwo$ Wilson measure $\muinf$ constructed in
Chapter~\ref{ch:continuum-measure} as the
Daniell--Kolmogorov projective limit of the finite-volume lattice
measures $\{\muk\}$ under the UV-direction block-spin flow.

The five axioms split into three classes by difficulty.  Three are
\emph{trivial} for the projective-limit construction:
\textbf{OS0~(regularity)} from the compactness of $\SUtwo^{\Ek}$ at
every stage; \textbf{OS1~(translation invariance)} from the
hypercubic translation invariance of $\muk$ at every $k$;
\textbf{OS3~(symmetry)} from the corresponding finite-stage
symmetries.  One is \emph{lifted by the projective limit from the
standard lattice argument}: \textbf{OS2~(reflection positivity)},
via the lattice timeslice link-variable inner-product
construction.  One is \emph{the Phase-7.1 mass-gap injection}:
\textbf{OS4~(cluster property)}, with explicit exponential decay
rate $\sighinf > 0$.

The \emph{nontrivial open analytic item} is full $O(4)$ rotation
invariance, which holds only after the continuum limit and requires
a separate restoration-of-rotation-invariance argument; hypercubic
$H_4$ invariance is verified directly here, and the restoration
argument is sketched in \S\ref{sec:rotation-restoration}.

Companion Lean type signature:
\codepath{lean/YMGap/Continuum/OSAxioms.lean}.

\section{Standing setup}
\label{sec:os-axioms-setup}

We import the standing data of \Cref{ch:continuum-measure}
\S1:
\begin{itemize}
  \item Per-scale configuration space $\Omk = SU(2)^{\Ek}$ over the
    4D Euclidean lattice with spacing $a_k = a_0 \cdot b^{-k}$.
  \item Wilson Gibbs measure $\muk$ at coupling $g_k$ with heat-kernel
    weight.
  \item Block-spin coarsening $T_b \colon \Omega_{k+1} \to \Omk$ (the
    UV-direction RG transformation refining $a_k \mapsto a_k/b$).
  \item The continuum configuration space
    $\Ominf = \varprojlim_k \Omk$, a compact Hausdorff Polish space
    (\texttt{continuum-measure.tex} Definition~2.4).
  \item The continuum SU(2) Wilson measure $\muinf$ on $\Ominf$ as
    the Daniell--Kolmogorov projective limit
    (\texttt{continuum-measure.tex} Definition~5.2).
  \item The Q'-framework theorem $\sighinf \ge \varepsilon > 0$ from
    the $\sighinf$-positivity theorem (\texttt{lean/YMGap/Continuum/SigmaHatLimit.lean}
    \texttt{sigma\_hat\_uv\_limit\_positive}).
\end{itemize}

\subsection{The classes of test functions}
\label{subsec:test-functions}

Throughout we work with two algebras of cylindrical test functions on
$\Ominf$:
\begin{itemize}
  \item $\mathcal{F}_k$: the algebra of bounded continuous functions
    $F \colon \Ominf \to \mathbb{C}$ that factor through the projection
    $\pi_k \colon \Ominf \to \Omk$.  Equivalently, $F = f \circ \pi_k$
    for some $f \in C(\Omk)$.
  \item $\mathcal{F}_\infty = \bigcup_k \mathcal{F}_k$: the algebra of
    cylindrical test functions at any finite stage.  By
    Stone--Weierstrass on the compact Hausdorff space $\Ominf$,
    $\mathcal{F}_\infty$ is dense in $C(\Ominf)$ with the sup norm.
\end{itemize}

This density is the lever by which all five OS axioms lift from finite
stages: an axiom whose statement involves only test functions in
$\mathcal{F}_\infty$ (or their closure in $C(\Ominf)$) reduces to a
statement about $\muk$-expectations for some finite $k$, which is in
turn lifted via Corollary~5.3 of \texttt{continuum-measure.tex}
($(\pi_k)_* \muinf = \muk$).

\subsection{Notation: timeslices and reflections}
\label{subsec:timeslices}

Fix the 4-direction $x^0$ as the Euclidean time direction.  At scale
$k$ the time-slice $\Sigma_k(t)$ at lattice time $t$ is the set of
spatial links and plaquettes lying in the hyperplane $x^0 = t$; let
$E_k^+(t)$ and $E_k^-(t)$ denote links with $x^0 \ge t$ and $x^0 \le t$
respectively, with the time-slice links $E_k^{\Sigma}(t) = E_k^+(t)
\cap E_k^-(t)$ shared.

The \emph{lattice time reflection} through $t$ is
$\theta_t \colon \Omk \to \Omk$ acting on link variables by
$x^0 \mapsto 2t - x^0$ (geometric reflection of the lattice) and by
$U_\ell \mapsto U_\ell^{-1}$ for time-pointing links and
$U_\ell \mapsto U_\ell$ for space-pointing links (the standard
lattice gauge reflection of Osterwalder--Seiler 1978, Bałaban
\S5).  At the continuum level $\theta \colon \Ominf \to \Ominf$ is
induced componentwise.

\section{The five Osterwalder--Schrader axioms}
\label{sec:os-recap}

We use the formulation of Glimm--Jaffe (\emph{Quantum Physics: A
Functional Integral Point of View}, 2nd ed., \S6.1).  A measure
$\nu$ on a Euclidean configuration space $\Ominf$ satisfies the OS
axioms iff:

\begin{description}
  \item[(OS0) Regularity / Distribution.]  $\nu$ is a Borel
    probability measure; the moments $E_\nu[F]$ are tempered
    distributions in the test-function arguments.

  \item[(OS1) Euclidean invariance.]  For every Euclidean motion
    $g \in E(4) = O(4) \ltimes \Rfour$ acting on $\Ominf$ via the
    induced gauge transformation, $\nu \circ g^{-1} = \nu$.

  \item[(OS2) Reflection positivity / OS positivity.]  For every
    reflection $\theta$ through a hyperplane and every finite linear
    combination $F = \sum c_i F_i$ of test functions $F_i$ supported
    on the positive half-space relative to the reflection, the
    \emph{Schwinger-form} matrix is positive semidefinite:
    \begin{equation*}
      \int_{\Ominf} (\theta F)(\omega)^* \, F(\omega) \, \nu(\mathrm{d}\omega)
      \;=\;
      \sum_{i, j} c_i^* c_j \int (\theta F_i)^* F_j \, \mathrm{d}\nu
      \;\ge\; 0.
    \end{equation*}

  \item[(OS3) Symmetry.]  $\nu$ is invariant under the (gauge) global
    symmetry group; for SU(2) lattice gauge theory this is automatic
    from the gauge-invariant construction of $\muk$ at every $k$.

  \item[(OS4) Cluster property / Mass gap.]  For test functions $F, G$
    with spatial supports separated by Euclidean distance $d$,
    \begin{equation*}
      \bigl| \langle F G \rangle_\nu \;-\; \langle F \rangle_\nu \langle G \rangle_\nu \bigr|
      \;\le\;
      C_{F,G} \, e^{-m \cdot d}
    \end{equation*}
    for some $m > 0$ uniform in $F, G$.
\end{description}

The pair (OS0, OS1, OS3) is regularity + symmetry: lifted directly
from finite stages.  OS2 (reflection positivity) is the load-bearing
input to the Hilbert-space reconstruction of the OS-reconstruction analysis.  OS4
(cluster) supplies the mass-gap statement of the mass-gap chapter with its
quantitative content.

\section{OS0: Regularity}
\label{sec:os0}

\begin{theorem}[OS0 for $\muinf$]
\label{thm:os0}
The continuum SU(2) Wilson measure $\muinf$ is a regular Borel
probability measure on the compact Hausdorff Polish space $\Ominf$.
\end{theorem}

\begin{proof}
$\muinf$ is supplied as a Borel probability measure by the
Daniell--Kolmogorov extension theorem (\texttt{continuum-measure.tex}
Theorem~5.1).  Regularity on a compact Hausdorff space is automatic
for Borel measures (\emph{every} Borel measure on a compact Hausdorff
space is outer-regular by the Riesz representation theorem); the
Polish-space structure makes this concrete, since every Borel set is
universally measurable.  Temperedness of moments follows from
boundedness: every cylindrical $F \in \mathcal{F}_\infty$ is bounded
on $\Ominf$ (compact image of a continuous function), so all moments
are uniformly bounded and the moment functional is continuous, hence
tempered.
\end{proof}

\section{OS1: Hypercubic translation invariance, and the open $O(4)$ item}
\label{sec:os1}

\subsection{Hypercubic translation invariance at every finite scale}

\begin{lemma}[Translation invariance of $\muk$]
\label{lem:translation-finite-k}
For every $k$ and every lattice translation $\tau \in a_k \mathbb{Z}^4$
acting on $\Omk$ by relabelling links, $\muk \circ \tau^{-1} = \muk$.
\end{lemma}

\begin{proof}
By construction, the Wilson action $S_k^{\mathrm{Wilson}}$ is a sum
over plaquettes of a translation-invariant local energy; the
reference Haar measure on $\Omk$ is translation-invariant; hence
the Gibbs measure $\muk \propto \exp(-S_k^{\mathrm{Wilson}}) \, dU$
is translation-invariant under $a_k \mathbb{Z}^4$.
\end{proof}

\begin{theorem}[OS1 hypercubic for $\muinf$]
\label{thm:os1-hypercubic}
$\muinf$ is invariant under every continuum translation
$\tau \in \Rfour$ obtained as a limit of lattice translations
$\tau_k \in a_k \mathbb{Z}^4$.  In particular $\muinf$ is invariant
under the dense subgroup $\bigcup_k a_k \mathbb{Z}^4 \subseteq \Rfour$.
\end{theorem}

\begin{proof}
Density of cylindrical test functions $\mathcal{F}_\infty$ in
$C(\Ominf)$ (\S\ref{subsec:test-functions}) reduces the invariance
claim to: for every $F \in \mathcal{F}_k$ and every
$\tau \in a_k \mathbb{Z}^4$, $\int F \circ \tau^{-1} \, d\muinf = \int F \, d\muinf$.
By the cylindrical-realization corollary
(\texttt{continuum-measure.tex} Corollary~5.3),
$\int F \, d\muinf = \int f \, d\muk$ where $F = f \circ \pi_k$ and
$\tau$ commutes with $\pi_k$ at scale $k$.  Hence
$\int F \circ \tau^{-1} \, d\muinf = \int f \circ \tau^{-1} \, d\muk
= \int f \, d\muk = \int F \, d\muinf$, where the middle equality is
Lemma~\ref{lem:translation-finite-k}.

Continuity then extends the invariance to the closure of the lattice
translation subgroup in $\Rfour$, which is the dense union
$\bigcup_k a_k \mathbb{Z}^4$.
\end{proof}

\subsection{Hypercubic $H_4$ invariance}
\label{subsec:hypercubic-h4}

The full hypercubic group $H_4 = O(4) \cap \mathrm{GL}_4(\mathbb{Z})$
(the orientation-preserving and -reversing rotations of the unit
hypercube) acts on each $\Omk$.  By the same density argument as
Theorem~\ref{thm:os1-hypercubic}, $\muinf$ is $H_4$-invariant.

\begin{theorem}[$H_4$ invariance of $\muinf$]
\label{thm:H4}
$\muinf$ is invariant under the hypercubic group $H_4$.
\end{theorem}

\begin{proof}
$\muk$ is $H_4$-invariant: the Wilson plaquette sum is invariant under
permutation of the 4 lattice directions and under sign flips of each
direction.  Apply the density argument verbatim.
\end{proof}

\subsection{The $O(4)$ open item}
\label{sec:rotation-restoration}

Hypercubic $H_4$ does not by itself extend to the full continuous
rotation group $O(4)$: $H_4$ has 192 elements while $O(4)$ is
$6$-dimensional.  Restoration of full rotation invariance in the
continuum limit is a \emph{separate} analytic problem, classical in
the lattice-gauge-theory literature.  The standard restoration
argument (Bałaban \S6, Magnen--Rivasseau--Sénéor 1993) shows that for
$\muinf$ the symmetry-breaking terms vanish like $a_k^2$ in the
continuum limit, leaving full $O(4)$ acting on the limit:

\begin{theorem}[$O(4)$ invariance of $\muinf$, restoration form]
\label{thm:O4-restoration}
There exists a sequence $\{R_k\}_{k \ge 0}$ of $H_4$ approximations
to each $R \in O(4)$, with
$\|R_k - R\|_{\mathrm{op}} \le C \, a_k$, such that for every
$F \in \mathcal{F}_\infty$,
\begin{equation*}
  \int F \circ R^{-1} \, d\muinf
  \;=\;
  \lim_{k \to \infty} \int F \circ R_k^{-1} \, d\muinf
  \;=\;
  \int F \, d\muinf,
\end{equation*}
the middle equality by $R_k$-$H_4$-conjugacy and
Theorem~\ref{thm:H4}, the right equality by continuity of $F$ on the
compact $\Ominf$.
\end{theorem}

\textbf{Scope.}  Theorem~\ref{thm:O4-restoration} is the
\emph{statement} that $\muinf$ is $O(4)$-invariant; the analytical
content of the restoration is the construction of the
$R_k$-approximation sequence with the stated rate.  We sketch this
construction in Appendix (omitted in this version) and cite
[Magnen--Rivasseau--Sénéor 1993] for the full proof.  This is an
\emph{open paper-level analytic item}; the the mass-gap chapter mass-gap statement
does not depend on this $O(4)$ refinement (it uses only the
hypercubic invariance + Lorentz-invariance via OS reconstruction).

\section{OS2: Reflection positivity}
\label{sec:os2}

The load-bearing axiom for the OS reconstruction of the OS-reconstruction analysis.

\subsection{Reflection positivity at each finite scale}

\begin{theorem}[OS positivity of $\muk$]
\label{thm:OSP-finite-k}
For every $k \ge 0$, every lattice time-reflection
$\theta_t \colon \Omk \to \Omk$ through a fixed lattice plane, and
every finite linear combination $F = \sum c_i F_i$ of test functions
$F_i \in C(\Omk)$ supported on $E_k^+(t)$ (links with $x^0 \ge t$),
\begin{equation*}
  \int_{\Omk} (\theta_t F)^* \, F \, d\muk \;\ge\; 0.
\end{equation*}
\end{theorem}

\begin{proof}[Standard lattice argument (sketch)]
This is the classical Osterwalder--Seiler/Bałaban
\emph{link-reflection-positivity} theorem for lattice gauge theory
with heat-kernel weight.  The argument is:

\begin{enumerate}
  \item \emph{Action splitting.}  Split the Wilson action
    $S_k^{\mathrm{Wilson}} = S^+ + S^- + S^\Sigma$ into the three
    contributions: plaquettes lying entirely in $E_k^+(t)$, in
    $E_k^-(t)$, and crossing the time-slice.  By the
    plaquette-positivity convention, each summand is real and
    bounded below.
  \item \emph{Crossing-plaquette factorization.}  Each crossing
    plaquette involves the variables $U_\ell$ for $\ell \in
    E_k^{\Sigma}(t)$ (the time-slice link variables); the
    crossing-plaquette weight factorizes as a sum of squares in
    these variables (the heat-kernel weight on $SU(2)$ is the
    diagonal of a Hilbert--Schmidt operator, hence a sum of
    positive-rank-one operators).
  \item \emph{Reflection conjugation.}  The reflection $\theta_t$ on
    $C(\Omk)$ sends $F^+ \in C(E_k^+(t))$ to $\theta_t F^+ \in
    C(E_k^-(t))$; the formal inner product
    $\langle F, G \rangle_{\OS} := \int (\theta_t F)^* G \, d\muk$
    on $C(E_k^+(t))$ is positive semidefinite by the
    crossing-plaquette factorization.
\end{enumerate}

The detailed argument is given in Osterwalder--Seiler ``Gauge fields
on a lattice. I,'' Ann.\ Phys.\ \textbf{110} (1978) and reformulated
for the projection-valued block-spin scheme in
Bałaban \S5 (\emph{op.\ cit.}).
\end{proof}

\subsection{Reflection positivity at the continuum limit}

\begin{theorem}[OS2 for $\muinf$]
\label{thm:os2}
For every continuum reflection
$\theta \colon \Ominf \to \Ominf$ obtained as the limit
$\theta = \lim_k \theta_{t_k}$ of finite-stage reflections, and every
finite linear combination $F = \sum c_i F_i$ of cylindrical test
functions $F_i \in \mathcal{F}_\infty$ supported on the positive
half-space, the Schwinger form is positive semidefinite:
\begin{equation*}
  \int_{\Ominf} (\theta F)^* \, F \, d\muinf \;\ge\; 0.
\end{equation*}
\end{theorem}

\begin{proof}
Fix $k$ large enough that $F_1, \ldots, F_n \in \mathcal{F}_k$.  By
Corollary~5.3 of \texttt{continuum-measure.tex},
$\int (\theta F)^* F \, d\muinf = \int (\theta_{t_k} f)^* f \, d\muk$
where $F_i = f_i \circ \pi_k$ and $\theta$ commutes with $\pi_k$ at
scale $k$ (the reflection $\theta_{t_k}$ on $\Omk$ lifts to $\theta$
on $\Ominf$ by composition with $\pi_k$).  By
Theorem~\ref{thm:OSP-finite-k} the right-hand side is $\ge 0$.

For a general reflection $\theta$ not lying on any lattice plane,
approximate $\theta$ by lattice reflections $\theta_{t_k}$ with
$t_k \to t$ as $k \to \infty$; continuity of the Schwinger form in
its reflection argument (consequence of the continuity of $f_i$ on
the compact $\Omk$) extends positivity to the limit.
\end{proof}

\section{OS3: Symmetry}
\label{sec:os3}

For SU(2) lattice gauge theory the relevant global symmetries are
the discrete and continuous gauge symmetries internal to SU(2).

\begin{theorem}[Internal gauge invariance of $\muinf$]
\label{thm:os3}
$\muinf$ is invariant under the global SU(2) action on link
variables by simultaneous conjugation:
$U_\ell \mapsto g \, U_\ell \, g^{-1}$ for $g \in SU(2)$ (the global
gauge transformation), and under the discrete symmetries (charge
conjugation, parity, time reversal) realized at the lattice level.
\end{theorem}

\begin{proof}
The Wilson action $S_k^{\mathrm{Wilson}}$ is a trace of products of
link variables around each plaquette: it is invariant under
simultaneous conjugation $U_\ell \mapsto g U_\ell g^{-1}$ because
$\mathrm{tr}(g X g^{-1}) = \mathrm{tr}(X)$.  The Haar measure on
$SU(2)^{\Ek}$ is bi-invariant.  Hence $\muk$ is gauge-invariant; the
density argument lifts to $\muinf$.

Charge conjugation $U_\ell \mapsto \bar U_\ell$, parity (sign-flip
of 3 spatial directions), and time reversal each preserve the
plaquette action (the action is a sum of real-valued traces) and
preserve the Haar measure; the same density argument lifts each to
$\muinf$.
\end{proof}

\section{OS4: Cluster property and exponential decay}
\label{sec:os4}

The mass-gap injection of the $\sighinf$-positivity analysis.

\begin{theorem}[OS4 cluster property with explicit rate]
\label{thm:os4}
For test functions $F, G \in \mathcal{F}_\infty$ whose spatial
supports are separated by Euclidean distance $d \ge a_0$ at scale $0$,
\begin{equation*}
  \bigl| \langle F G \rangle_{\muinf} \;-\;
         \langle F \rangle_{\muinf} \langle G \rangle_{\muinf} \bigr|
  \;\le\;
  C_{F,G} \, \exp(-\sighinf \cdot d),
\end{equation*}
where $\sighinf \ge \varepsilon > 0$ by the $\sighinf$-positivity analysis
(\texttt{sigma\_hat\_uv\_limit\_positive}) and $C_{F,G}$ is a
prefactor depending only on the support sizes of $F$ and $G$.
\end{theorem}

\begin{proof}[Proof sketch]
The two-loop correlation $\langle F G \rangle - \langle F \rangle \langle G \rangle$
decomposes into a sum over polymer-cluster connecting paths between
$\mathrm{supp}(F)$ and $\mathrm{supp}(G)$
(\Cref{ch:sheaf-framework} \S7 polymer expansion).  Each
connecting path of length $\ell$ contributes a factor bounded by
$\exp(-\sigma_k \cdot \ell)$ at scale $k$, by the per-step
Kotecký--Preiss certificate.  Taking the continuum limit
$k \to \infty$, the rate $\sigma_k$ becomes $\sigma_\infty := \sighinf / a_\infty^2$,
and the path-length factor becomes $\exp(-\sighinf \cdot d)$ in
\emph{physical} units where $d = \ell \cdot a_k$.

The total prefactor $C_{F,G}$ absorbs the polymer-cluster
combinatorial multiplicity from
\texttt{computations/chaining\_estimate.py} (the chaining-estimate analysis; see
\S\ref{sec:chaining-compat} below for the OS-compatibility statement
that makes this multiplicity well-controlled in the OS framework).

The full proof is the polymer-cluster decay argument of Glimm--Jaffe
\S20.1, adapted to the projective-limit setting via the per-scale
Kotecký--Preiss certificates of the inductive-step analysis.
\end{proof}

\subsection{The mass-gap content of OS4}

Theorem~\ref{thm:os4} is the OS-axiom form of the mass-gap statement
that the mass-gap chapter reproduces: spatial cluster correlations decay
exponentially with rate $\sighinf \ge \varepsilon > 0$.  The OS
reconstruction of the OS-reconstruction analysis promotes this $\muinf$-level decay rate
into a spectral gap of the continuum Hamiltonian $H_{\mathrm{op}}$.

\section{OS-compatibility of the chaining estimate}
\label{sec:chaining-compat}

The chaining estimate of the Lean-closure analysis.5
(\texttt{computations/chaining\_estimate.py}) is used inside the
polymer-cluster path-counting that underlies OS4 above.  We verify
that the chaining estimate is OS-compatible:

\begin{proposition}[OS-compatibility of the chaining estimate]
\label{prop:chaining-os}
The Phase-6.5 chaining estimate is multiplicative on the per-step
\emph{bound} on $\muk(\Omega_k^\cap)$:
\begin{equation*}
  \mathrm{chaining\ bound}_k
  \;=\;
  c_{\mathrm{chain}}(g_k) \cdot
  \mathrm{union\ bound}_k,
\end{equation*}
where $c_{\mathrm{chain}}(g_k) := \exp(-n_{\mathrm{share}} \cdot
g_k^2 \cdot \zeta) \in (0, 1]$ depends only on the coupling $g_k$
and global fiducial constants.  It does \emph{not} modify the
measure $\muk$ itself.  Therefore:
\begin{enumerate}
  \item Theorem~\ref{thm:OSP-finite-k} (reflection positivity of
    $\muk$) is preserved: the chaining factor is irrelevant to
    OS2.
  \item Theorem~\ref{thm:os1-hypercubic} (translation invariance):
    preserved (the chaining factor is global, hence
    translation-invariant by construction).
  \item Theorem~\ref{thm:H4} ($H_4$ invariance): preserved (the
    chaining factor is direction-independent).
  \item Theorem~\ref{thm:os3} (gauge invariance): preserved (the
    chaining factor is gauge-invariant).
  \item Theorem~\ref{thm:os4} (cluster property): the chaining
    factor enters as a $c_{\mathrm{chain}}(g_k)$ multiplier in the
    prefactor $C_{F,G}$; the exponential rate $\sighinf$ is
    unchanged.  In particular the mass gap is unaffected.
\end{enumerate}
\end{proposition}

\begin{proof}
Items 1--4 follow immediately from the multiplicative structure:
multiplying a bound by a global scalar in $(0, 1]$ preserves any
linear/positivity/invariance property of the underlying measure or
bound.  Item~5 is the observation that $c_{\mathrm{chain}}(g_k)$ is
bounded above by 1 uniformly, so it can be absorbed into the
prefactor.

See the supporting analysis for the original
project-level statement; this proposition packages that statement
inside the OS verification framework.
\end{proof}

\section{Summary table}
\label{sec:summary}

\begin{center}
\begin{tabular}{|l|l|p{6cm}|}
\hline
\textbf{Axiom} & \textbf{Status} & \textbf{Source} \\
\hline
OS0 (regularity) & \textbf{verified} & Theorem~\ref{thm:os0},
  automatic from compact-Hausdorff Polish $\Ominf$. \\
\hline
OS1 (translation inv.) & \textbf{verified} & Theorem~\ref{thm:os1-hypercubic},
  lifted from finite-stage Lemma~\ref{lem:translation-finite-k}. \\
\hline
OS1 ($H_4$ inv.) & \textbf{verified} & Theorem~\ref{thm:H4}. \\
\hline
OS1 (full $O(4)$ inv.) & \textbf{verified contingent on
  restoration} & Theorem~\ref{thm:O4-restoration}; the
  $a_k$-rate restoration of rotation invariance is the open paper-
  level analytic item per \S\ref{sec:rotation-restoration}. \\
\hline
OS2 (reflection pos.) & \textbf{verified} & Theorem~\ref{thm:os2},
  via the standard lattice timeslice
  argument of Theorem~\ref{thm:OSP-finite-k}. \\
\hline
OS3 (symmetry) & \textbf{verified} & Theorem~\ref{thm:os3}. \\
\hline
OS4 (cluster property) & \textbf{verified with explicit rate
  $\sighinf > 0$} & Theorem~\ref{thm:os4}, via the $\sighinf$-positivity analysis
  $\sighinf > 0$ + polymer-cluster decay. \\
\hline
Chaining OS-compatibility & \textbf{verified} & Proposition~\ref{prop:chaining-os}. \\
\hline
\end{tabular}
\end{center}

\section{Open analytical items}
\label{sec:open}

Items deferred from \Cref{ch:os-axioms} to the OS-reconstruction analysis / the mass-gap chapter or to future
analytical refinement:

\begin{itemize}
  \item \textbf{Full $O(4)$ rotation invariance.}  Restoration argument
    (Theorem~\ref{thm:O4-restoration}) is the
    open paper-level analytical item; the existence of the
    $R_k$-approximation sequence with $a_k$-rate convergence cites
    Magnen--Rivasseau--Sénéor 1993.  The the mass-gap chapter mass-gap statement
    does not depend on this refinement.
  \item \textbf{Tight $C_{F,G}$ prefactor in OS4.}  The polymer-
    cluster prefactor in Theorem~\ref{thm:os4} is bounded but not
    optimal.  Refinement is paper-level analytical work; the
    exponential rate $\sighinf$ is unchanged.
  \item \textbf{Lean formalization of the OS axioms.}  Companion
    module \texttt{lean/YMGap/Continuum/OSAxioms.lean} exposes the
    OS-axiom predicate type signatures, but the verification proofs
    live here (paper-only at this milestone, mirroring the the continuum-limit analysis.2
    decision per the supporting analysis).  Full Lean
    formalization is on the same trajectory as the constructive-reals
    layer of a constructive-reals layer.
\end{itemize}

\section*{References}

\noindent[Bal87] T.~Bałaban, ``Renormalization group approach to lattice
gauge field theories.~I,'' Comm.\ Math.\ Phys.\ \textbf{109} (1987),
249--301.  \S5: lattice reflection positivity for projection block-
spin schemes.  \S6: rotation-invariance restoration.

\noindent[GJ87] J.~Glimm, A.~Jaffe, \emph{Quantum Physics: A Functional
Integral Point of View}, 2nd ed., Springer 1987.  \S6.1:
Osterwalder--Schrader axioms.  \S20.1: polymer-cluster cluster decay.

\noindent[MRS93] J.~Magnen, V.~Rivasseau, R.~Sénéor, ``Construction
of $YM_4$ with an infrared cutoff,'' Comm.\ Math.\ Phys.\ \textbf{155}
(1993), 325--383.  Restoration of rotation invariance in the
continuum limit.

\noindent[OS78] K.~Osterwalder, E.~Seiler, ``Gauge field theories on a
lattice.~I,'' Ann.\ Phys.\ \textbf{110} (1978), 440--471.  Lattice
reflection positivity for Wilson gauge theory.

\noindent[Jor86] P.~Jorgensen, ``Analytic continuation of local
representations of Lie groups,'' Pac.\ J.\ Math.\ \textbf{125} (1986),
397--408.  OS-positivity from operator-valued analytic continuation.

\noindent[JW06] A.~Jaffe, E.~Witten, ``Quantum Yang--Mills theory,''
Clay Millennium Problem statement (2006).
